\section{Follow-up: where the 7B result ends}
\label{app:followup}

Section~\ref{sec:discussion} names three things that would change our mind. The
first was ``a GRPO run at this scale with an order of magnitude more prompt
coverage, reaching a plateau in its own reward curve, that beats
$\max\mathcal{B}$ at any operating point.'' We ran it. It does.

This appendix reports that result and a second run that isolates why. Both are
exploratory: they use the frozen test set at an operating point the protocol
already spent, they were specified after the 17 registered groups were complete,
and they do not enter $\Delta$ or any count in the main text. We report them
because we said we would, and because the answer bounds the paper's claim in a
way a reader is entitled to know.

\subsection{Design}

We retested the single cell with the \emph{smallest} registered margin:
Config~3 at $m{=}64$, seed~1, where $\max\mathcal{B}$ is continued SFT at
$80.50\%$ and the registered GRPO run scored $79.43\%$ for $\Delta = -1.06$.
This is the cell most favourable to GRPO and therefore the sharpest test: if
more compute cannot flip it here, it will not flip the others.

Three runs, all branching from the same warm start $W_{64}$ on the same prompt
pool, differing only in budget and in which prompts they see:

\begin{itemize}
\item \textbf{Two runs at ten times the wall clock} (57{,}600\,s against
5700\,s), with two different training seeds and the demonstration seed held
fixed, so $W_m$ and the pool are identical to the registered run.
\item \textbf{One run at the registered budget with a difficulty-filtered
pool.} Section~\ref{sec:analysis} shows that $73.7\%$ of this cell's prompt
groups carry zero advantage. We scored the pool once with $W_{64}$ under greedy
decoding, kept only the $440$ of $2382$ prompts it answered incorrectly, and
trained on that subset under the \emph{same} 5700\,s cap. This asks whether the
registered result is an artifact of wasted gradient signal rather than of
budget.
\end{itemize}

% Generated by analysis/gen_followup_table.py. Do not edit by hand.
\begin{table}[t]
\centering
\small
\setlength{\tabcolsep}{4pt}
\caption{Follow-up runs on the Config~3 cell with the smallest registered margin ($m{=}64$, seed~1), where $\max\mathcal{B}$ is continued SFT at $80.50\%$. All three are exploratory and do not enter the pre-registered decision quantity. Coverage is prompt groups divided by the 2382-item pool.}
\label{tab:followup}
\begin{tabular}{lrrrrrrr}
\toprule
run & steps & wall (s) & coverage & zero-var & informative & test acc & $\Delta$ \\
\midrule
registered & 39 & 5750 & 0.13 & 73.7\% & 82 & 79.43 & $-1.06$ \\
\addlinespace
10$\times$ coverage, seed A & 414 & 57667 & 1.39 & 86.0\% & 462 & 87.23 & $\mathbf{+6.74}$ \\
10$\times$ coverage, seed B & 415 & 57694 & 1.39 & 86.1\% & 461 & 86.88 & $\mathbf{+6.38}$ \\
difficulty-filtered, matched budget & 33 & 5753 & 0.11 & 54.5\% & 120 & 81.56 & $\mathbf{+1.06}$ \\
\bottomrule
\end{tabular}
\end{table}


\subsection{What happened}

\textbf{At ten times the budget, GRPO wins, decisively and reproducibly.}
Both seeds reach $87.23\%$ and $86.88\%$ against continued SFT's $80.50\%$:
$\Delta = +6.74$ and $+6.38$, McNemar $p = 0.0013$ and $0.0029$, bootstrap
intervals $[+2.84, +10.64]$ and $[+2.48, +10.28]$, and a bootstrap probability
of clearing the $+3$ Go threshold of $0.97$ and $0.95$. By the paper's own
pre-registered criterion this is a Go. Coverage rose from $0.13$ epochs to
$1.39$.

\textbf{At the registered budget, fixing the degeneracy changes nothing
measurable.} The difficulty filter worked as intended: zero-advantage groups
fell from $73.7\%$ to $54.5\%$ and informative prompts rose from 82 to 120,
while the run consumed \emph{less} on every compute axis than the registered one
(33 steps against 39, 2112 verifier calls against 2496, 5753\,s against
5750\,s). Accuracy rose to $81.56\%$, nominally $+1.06$ over continued SFT. But
McNemar gives $p = 0.678$ and the interval is $[-2.13, +4.26]$: this is
indistinguishable from the registered $-1.06$, and $P(\Delta \geq +3) = 0.12$.
The sign moved; nothing else did.

\textbf{More compute did not make the RL arm more efficient.} The
zero-advantage fraction \emph{rose} with training, from $73.7\%$ to $86.0\%$, as
the policy solved more of the pool outright. Informative prompts grew only
$5.6\times$ across a $10.6\times$ increase in steps. The win at ten times the
budget was bought by brute force, not by better use of the signal.

\subsection{What this does to the paper}

It bounds the 7B claim and leaves the rest standing.

\emph{The registered finding is unchanged as registered.} Under the
pre-registered protocol---equal wall clock for every arm---GRPO does not clear
$+3$ at 7B, and the compute-matched follow-up confirms that this is not an
artifact of the degeneracy we diagnosed. Both the original $-1.06$ and the
filtered $+1.06$ sit inside the same interval.

\emph{The 7B configuration no longer supports a claim about scale.} It was
included to answer ``your models are too small,'' and it cannot: at $10\times$
the budget the same cell reverses to a decisive Go. What Config~3 demonstrates
is a \emph{budget} boundary, not a capability one, and
Section~\ref{sec:scale} now says so. The quantitative core of this paper is the
11 groups at 0.5B and 1B, where $\max\mathcal{B}$ is attained by a deterministic
arm leading by $8.5$ to $36.9$ points rather than by $1$ to $4$, and where GRPO
had $0.53$ to $0.79$ epochs rather than $0.13$.

\emph{The paper's own thesis absorbs this.} The claim is that a reported RLVR
effect depends on the comparison it is measured against. We can now put three
numbers on that dependence, from the same runs and the same test set: changing
which control arm is nominated moves the reported effect by $19.5$ points
(Section~\ref{sec:statistics}); changing the compute budget by an order of
magnitude moves it by $7.8$; and changing the inference engine alone moves it by
$3.0$ (Appendix~\ref{subsec:enginegate}). A reader who takes only one thing from
this paper should take that an RLVR number without its comparison, its budget
and its decoding stack attached is not yet a result.
